Next define a homomorphism @@M0@@ by assigning to @@M1@@ the element @@M2@@; this homomorphism vanishes on @@M3@@, and consequently defines a homomorphism @@M4@@. Finally, define a homomorphism @@M5@@ by assigning to @@M6@@ the element @@M7@@; this homomorphism vanishes on @@M8@@ and therefore defines a homomorphism @@M9@@.

\sgaStableTarget{sga6.e07.prop.4-7-6}\textbf{Proposition 4.7.6.}
\begin{enumerate}[label=\roman*)]
\item @@M10@@ \uline{and} @@M11@@ \uline{are} @@M12@@-\uline{homomorphisms};
\item @@M13@@ \uline{is} @@M14@@-\uline{linear and} @@M15@@ \uline{is} @@M16@@-\uline{linear};
\item @@M17@@ \uline{is the homomorphism} ``\uline{multiplication by} @@M18@@'';
\item \uline{for every} @@M19@@, \uline{and every} @@M20@@, @@M21@@;
\item @@M22@@.
\end{enumerate}

We need only set the machinery in motion.

Returning to the geometric situation recalled above, we see that @@M23@@ is constructed from @@M24@@, @@M25@@, the maps @@M26@@ and @@M27@@, and the element @@M28@@ of @@M29@@ by the preceding method; the @@M30@@-operations on @@M31@@, however, agree with the @@M32@@-operations defined above on @@M33@@ only up to torsion, as long as relation (4.3.1) is proved only modulo torsion.

\sgaStableTarget{sga6.e07.thm.4-8}\textbf{Theorem 4.8.} \uline{Let} @@M34@@ \uline{be a regular closed immersion of constant codimension} @@M35@@, @@M36@@ \uline{the scheme obtained by blowing up} @@M37@@, \uline{and} @@M38@@, \uline{so that we have the commutative diagram}:
@@M39@@
\uline{Suppose that} @@M40@@ \uline{is quasi-compact and has an ample sheaf}. \uline{Then for every} @@M41@@ \uline{we have the exact sequence}:
